%% ============================================================
%%  latest.tex
%%  "The Interest Parameter in Perpetual Futures:
%%   Shariah Analysis and Empirical Evidence from
%%   Centralized and Decentralized Exchanges"
%%  Ahmed, Bhuyan & Islam — February 2026
%% ============================================================
\documentclass[10pt,twocolumn]{article}
\usepackage[utf8]{inputenc}
\usepackage[margin=0.75in]{geometry}
\usepackage{amsmath,amssymb}
\usepackage{graphicx}
\usepackage{hyperref}
\usepackage[numbers]{natbib}
\usepackage{booktabs}
\usepackage{abstract}
\usepackage{xurl}
\usepackage{titlesec}
\usepackage{authblk}
\usepackage{xcolor}
\usepackage{array}
\usepackage[protrusion=true,expansion=false]{microtype}
\usepackage{enumitem}
\usepackage{caption}
\usepackage{fancyvrb}
\usepackage{multirow}

%% Section spacing
\titlespacing*{\section}{0pt}{13pt}{6pt}
\titlespacing*{\subsection}{0pt}{10pt}{4pt}
\titlespacing*{\subsubsection}{0pt}{8pt}{3pt}

%% Abstract style
\renewcommand{\abstractnamefont}{\normalfont\bfseries\large}
\renewcommand{\abstracttextfont}{\normalfont\small}

%% Math operators
\DeclareMathOperator{\clamp}{clamp}

\definecolor{ribared}{RGB}{180,30,30}
\definecolor{halalgreen}{RGB}{30,120,60}
\definecolor{warnambr}{RGB}{160,100,0}

\hypersetup{colorlinks=true,linkcolor=blue,citecolor=blue,urlcolor=blue}

%% ============================================================
\title{\LARGE\textbf{The Interest Parameter in Perpetual Futures:
Shariah Analysis and Empirical Evidence from
Centralized and Decentralized Exchanges}}

\author[1]{Shehzad Ahmed}
\author[2]{Rafiqul Bhuyan}
\author[3]{Rubaiyat Islam}

\affil[1]{Department of Finance, Independent University, Bangladesh\\
\small\texttt{2120922@iub.edu.bd}}
\affil[2]{Department of Finance, Alabama A\&M University, USA; Independent University, Bangladesh\\
\small\texttt{rafiqul.bhuyan@aamu.edu}}
\affil[3]{Department of Software Engineering, Daffodil International University,
Bangladesh\\ \small\texttt{islam.swe@diu.edu.bd}}

\date{February 2026}
%% ============================================================

\begin{document}

\twocolumn[\begin{@twocolumnfalse}
\maketitle\thispagestyle{empty}

\begin{abstract}
\noindent
We document a previously unrecognised structural bifurcation in the \$100
billion perpetual futures market that has decisive implications for Islamic
finance: the interest (\textit{riba}) component embedded in every
centralized-exchange (CEX) funding formula is neither mathematically necessary
nor universal. A recent formal proof demonstrates that spot-price convergence
is achievable with zero interest \cite{Ackerer2024}. Consistent with this
result, decentralized exchange (DEX) protocols dYdX~v4 and D8X set the
interest parameter to exactly zero --- a fact documented in their official
APIs and confirmed by 365 days of observed funding rates.

We establish a \textbf{four-category taxonomy}: \emph{CEX perpetuals} (Binance,
Bybit, OKX) embed $I = 0.01\%$ per 8 hours (= 10.95\%/year structural
\textit{riba}); \emph{Premium-only DEX perpetuals} (dYdX~v4, D8X) set $I = 0$
and are structurally \textit{riba}-free; \emph{Hybrid DEX perpetuals}
(Hyperliquid) use the identical CEX formula despite decentralised architecture
and are therefore \textit{riba}-bearing; and \emph{DEX expiring futures}
(MCDEX, Serum-based) carry no funding mechanism at all and combine a fixed
settlement date with smart-contract self-custody, representing the theoretically
strongest Shariah candidate. This taxonomy resolves the persistent error in the
Islamic finance literature, which has treated all leveraged crypto derivatives
as a monolithic instrument.

Empirical analysis of 40,218 funding intervals (February 2025 -- February 2026)
across five platforms produces a clean natural experiment: three CEX platforms
(Binance, Bybit, OKX) with $I > 0$ versus dYdX~v4 with $I = 0$ exhibit
statistically distinct distributions ($t = 59.95$, $p < 10^{-300}$, Cohen's
$d = 0.725$, KS $= 0.487$), while Hyperliquid --- a DEX with $I > 0$ ---
clusters with CEX rather than dYdX, confirming that \emph{formula structure},
not exchange architecture, determines the Shariah classification. A pivotal signal: SeekersGuidance (2025),
the most rigorous contemporary conservative ruling, explicitly concedes that
removing the interest component resolves the \textit{riba} objection --- thereby
validating our core empirical finding as jurisprudentially decisive.

Principles-based Shariah analysis confirms that premium-only DEX perpetuals,
when used on economically productive underlying assets with disciplined risk
management, satisfy El-Gamal's riba conditions, Kamali's standardisation
criterion, and Salamon's maysir test. This paper demonstrates that Shariah
compliance in perpetual futures is achievable through proactive structural
design --- specifically, the elimination of the fixed interest parameter from
the funding formula --- and provides a principled framework for evaluating,
structuring, and implementing compliant perpetual contracts.

\medskip
\noindent\textbf{Keywords:} \textit{riba}, \textit{gharar}, \textit{maysir},
perpetual futures, Islamic finance, DeFi, DEX, Shariah compliance, funding rate,
decentralized exchange, halal trading
\end{abstract}
\bigskip
\end{@twocolumnfalse}]

%% ============================================================
\section{Introduction}
%% ============================================================

Perpetual futures are the defining financial instrument of the cryptocurrency
era. Daily trading volumes exceed \$100 billion, dwarfing spot markets and
providing the principal mechanism for price discovery and leveraged exposure
\cite{Chainalysis2025,CoinGecko2025}. Yet for 1.8 billion Muslims worldwide
--- including a rapidly expanding cohort of quantitatively sophisticated
finance professionals --- this market has been categorically off-limits. The
prohibition is not rooted in speculation alone; it is rooted in a specific
technical feature: every major centralized perpetual exchange has embedded a
fixed, predetermined interest charge into the structural core of its funding
mechanism. That interest charge is \textit{riba}. Under every school of Islamic
jurisprudence, it is prohibited.

Islamic scholars reached this conclusion correctly. Binance labels the charge
\texttt{interestRate} in its public API. The value is always
\texttt{0.00010000} --- $0.01\%$ per 8 hours, $0.03\%$ per day, $10.95\%$ per
year --- regardless of market conditions. It is predetermined, fixed,
labeled as interest, and structurally unavoidable. The conservative consensus
\cite{Musaffa2024,IslamQA2023,PracticalIF2023,Indonesian2024,Alias2021} is
not a misreading of the instrument: it is a correct reading of the instrument
as it exists on CEX platforms.

This paper demonstrates that the interest component can be structurally
eliminated from the perpetual futures funding formula through deliberate
platform design. Premium-only DEX protocols --- most notably dYdX~v4 and D8X
--- set the interest parameter to exactly zero, producing a funding mechanism
that is market-driven, bilateral, and free of any fixed predetermined charge.
The rate is pure market premium: it can be positive, zero, or negative,
depending entirely on the relative price of the perpetual versus the spot index.
Formal mathematics confirms this is not a functional compromise: Ackerer,
Hugonnier \& Jermann \cite{Ackerer2024} proved in 2024 that spot convergence is
achievable with $I = 0$, establishing that the CEX interest floor was a
commercial design choice, not a technical requirement. The implication for
Islamic finance is direct: Shariah compliance is not a constraint that perpetual
futures cannot satisfy --- it is a design criterion that has already been
satisfied, and can be replicated by any platform that adopts the premium-only
formula.

\subsection*{The Central Error in Existing Scholarship}

A review of 38 sources --- encompassing \textit{fatwas}, academic papers,
textbooks, and institutional standards --- reveals a consistent pattern: no
prior Islamic finance work distinguishes between CEX and DEX funding mechanisms.
Every existing prohibition treats perpetual futures as a single instrument.
The instrument being analysed is, without exception, the CEX instrument. The
DEX alternative is invisible in the literature, despite being commercially live
since 2023. The error is one of \emph{extrapolation}: correct diagnosis of CEX
\textit{riba} followed by incorrect generalisation to all perpetuals.

This paper corrects that error. We do not challenge the consensus on CEX
perpetuals --- it is right. We demonstrate that the consensus does not extend
to the premium-only DEX implementations that now operate at scale.

\subsection*{Contributions}

This paper makes four primary contributions.

\begin{enumerate}[leftmargin=*, noitemsep]
\item \textbf{Four-category taxonomy.} We develop a principled taxonomy
classifying leveraged crypto derivatives along two independent Shariah-relevant
dimensions: the interest parameter $I$ in the funding formula, and the presence
or absence of a settlement date. The resulting four categories --- CEX perpetual,
Premium-only DEX perpetual, Hybrid DEX perpetual, and DEX expiring futures ---
provide a replicable framework for evaluating any current or future platform.

\item \textbf{Natural experiment.} We exploit the existence of two structurally
similar DEX platforms with different interest parameters (dYdX $I=0$;
Hyperliquid $I=0.01\%$/8h) as a natural experiment, allowing us to isolate the
effect of the interest parameter on distributional outcomes independently of
market conditions --- empirically validating that formula design, not exchange
architecture, determines the Shariah-relevant properties of the funding mechanism.

\item \textbf{365-day empirical validation.} We collect and analyse 40,218
funding rate intervals across five platforms over a complete annual cycle,
confirming that the structural difference between premium-only and
interest-bearing formulas is persistent across both bull and bear market regimes,
not an artefact of short-term conditions.

\item \textbf{Reform-oriented implementation framework.} We develop a structured
platform-selection, asset-selection, and risk-management framework that
translates the empirical findings into actionable guidance for Muslim traders,
institutions, and Shariah boards evaluating perpetual futures instruments.
\end{enumerate}

%% ============================================================
\section{How Perpetual Futures Work and Where Riba Enters}
\label{sec:mechanics}
%% ============================================================

\subsection{The Funding Mechanism}

A perpetual futures contract has no expiry date. Without scheduled settlement,
the contract price would drift arbitrarily from the underlying spot price. The
funding mechanism is the engineering solution: every 8 hours (on CEX) or every
hour (on most DEX), the party whose position is ``on the expensive side'' pays
the counterparty an amount proportional to the price gap. This continuous
payment keeps the perpetual price tethered to spot.

The decisive question for Shariah compliance is not whether a payment occurs ---
it is \emph{how the payment is calculated}.

\subsection{The CEX Formula and Its Riba}

Binance \cite{Binance2019}, Bybit, and OKX all implement the same formula,
originating from BitMEX's 2016 launch \cite{Hayes2016}:
\begin{equation}
F = P + \clamp(I - P,\;-0.05\%,\;+0.05\%)
\label{eq:cex}
\end{equation}
where $F$ is the funding rate, $P$ is the premium index (perpetual price minus
spot, divided by spot), and $I = 0.01\%$ per 8 hours is the \textbf{fixed
interest rate}. Binance's public API returns:

\begin{Verbatim}[fontsize=\small]
{
  "interestRate": "0.00010000",
  "lastFundingRate": "0.00005102"
}
\end{Verbatim}

The field is named \texttt{interestRate}. The value is always
\texttt{0.00010000}. This is not interpretation; it is documentation. Three
Shariah characteristics of $I$ make it unambiguous \textit{riba}
\cite{ElGamal2006,Uddin2015}:

\begin{itemize}[noitemsep, leftmargin=*]
\item \textbf{Predetermined}: $I$ is fixed at $0.01\%$/8h regardless of market
conditions, counterparty creditworthiness, or economic outcome.
\item \textbf{Independent}: $I$ operates identically whether the market rises,
falls, or is flat. No market outcome can eliminate it.
\item \textbf{Labeled}: Binance itself calls it \texttt{interestRate} in its
API. No scholarly inference is required.
\end{itemize}

Figure~\ref{fig:interest} illustrates the API evidence. The structural riba
burden is $I \times 3 \times 365 = 10.95\%$/year --- an unavoidable floor that
no trading strategy can escape within CEX architectures.

\begin{figure}[htbp]
\centering
\includegraphics[width=\columnwidth]{fig_interest_evidence.png}
\caption{Binance \texttt{/fapi/v1/premiumIndex} API response. The
\texttt{interestRate} field --- always \texttt{0.00010000} ($0.01\%$/8h) ---
is the documentary source of \textit{riba}. This is Binance's own designation,
not a Shariah inference. The structural annual burden is $0.01\%\times3\times365
= 10.95\%$/year.}
\label{fig:interest}
\end{figure}

\subsection{The Mathematical Proof That Riba Is Unnecessary}

A critical 2024 result from financial mathematics directly addresses whether $I
> 0$ is technically required. Ackerer, Hugonnier \& Jermann
\cite{Ackerer2024} prove rigorously --- in both discrete and continuous time ---
that a perpetual futures contract converges to the spot price when $I = 0$.
The premium-only formula $F = P$ is provably sufficient for market function.
This transforms the Shariah analysis: the CEX interest floor is not an
engineering necessity. It is a \emph{design choice} --- a commercial decision
made by Arthur Hayes at BitMEX in 2016 \cite{Hayes2016} to attract market
makers, not a technical constraint. Riba was built in; it did not have to be.

\begin{quote}
\textbf{Theorem (Ackerer, Hugonnier \& Jermann 2024).} In both discrete and
continuous time, with premium intensity $\kappa_t > 0$ and zero interest
parameter $\iota = 0$, the no-arbitrage perpetual futures price converges to
the spot price $x_t$ whenever $r_a = r_b$. Specifically, in the
constant-parameter continuous case:
\[
f_t = \frac{\kappa}{\kappa + r_b - r_a}\,x_t,
\]
which equals $x_t$ when $r_a = r_b$. The premium-only funding rule is
therefore a mathematically sufficient condition for spot-price convergence;
the fixed interest parameter $I > 0$ used by CEX platforms is a commercial
design choice with no analytical necessity \cite{Ackerer2024}.
\end{quote}

Table~\ref{tab:ackerer} summarises the convergence conditions across all
three contract types covered by Ackerer et al.\ With $\iota = 0$, each
formula reduces to a ratio involving only $\kappa$ and the rate differential
$(r_b - r_a)$; when rates are aligned the ratio collapses to unity and
$f_t = x_t$ exactly. Our empirical finding that dYdX~v4 mean funding is
$0.00077\%$/h is consistent with Proposition~3 for linear contracts:
stablecoin-margined perpetuals imply $r_a \approx r_b$, pushing the ratio
toward one.

\begin{table}[htbp]
\centering
\caption{Ackerer et al.\ (2024) $\iota = 0$ convergence conditions by
contract type (continuous-time, constant parameters). Premium intensity
$\kappa$ alone anchors the futures price to spot when interest rates align.}
\label{tab:ackerer}
\small
\setlength{\tabcolsep}{4pt}
\renewcommand{\arraystretch}{1.6}
\begin{tabular}{@{}p{1.2cm}p{3.3cm}p{2.2cm}@{}}
\toprule
\textbf{Type} & \textbf{Price ($\iota=0$, constants)} & \textbf{Convergence} \\
\midrule
Linear &
  $\dfrac{\kappa}{\kappa + r_b - r_a}\,x_t$ &
  $r_a = r_b$ \\[4pt]
Inverse &
  $\dfrac{\kappa_I + r_a - r_b}{\kappa_I}\,x_t$ &
  $r_a = r_b$ \\[4pt]
Quanto &
  $\dfrac{\kappa}{\kappa + r_c - \sigma_x\sigma_z - r_a}\,z_t$ &
  $r_c = r_a {+} \sigma_x\sigma_z$ \\
\bottomrule
\multicolumn{3}{@{}l@{}}{\scriptsize In all cases $f_t = x_t$ (or $z_t$) when the convergence condition holds.}\\
\end{tabular}
\renewcommand{\arraystretch}{1}
\end{table}

\subsection{The DEX Formula: Premium Only}

dYdX~v4 and D8X implement the degenerate case of Equation~\eqref{eq:cex} with
$I = 0$:
\begin{equation}
F = P = \frac{P_{\mathrm{perp}} - P_{\mathrm{spot}}}{P_{\mathrm{spot}}}
\label{eq:dex}
\end{equation}
The funding rate is purely the signed market premium. When the perpetual trades
above spot (bullish sentiment), longs pay shorts; when it trades below spot
(bearish), shorts pay longs. The payment is bilateral, conditional on market
conditions, and can be zero or negative --- none of the three \textit{riba}
characteristics are present. The dYdX API returns no \texttt{interestRate}
field; there is none to return \cite{dYdX2024}.

\subsection{Platform Taxonomy}

Not all perpetual futures platforms are structurally equivalent, and not all
derivatives on decentralised protocols are perpetuals. Our survey yields a
four-category taxonomy (Table~\ref{tab:taxonomy}) that classifies instruments
by two independent dimensions: the presence of a fixed interest parameter $I$
in the funding formula, and whether a settlement date exists.

\begin{table}[htbp]
\centering
\caption{Four-category platform taxonomy. Shariah classification is determined
by the interest parameter $I$ and the presence of a settlement date, not by
CEX/DEX architecture alone.}
\label{tab:taxonomy}
\footnotesize
\setlength{\tabcolsep}{2pt}
\begin{tabular}{@{}p{1.45cm}p{1.6cm}p{1.4cm}p{1.05cm}p{1.35cm}@{}}
\toprule
\textbf{Category} & \textbf{Platforms} & \textbf{Riba source} & \textbf{Expiry} & \textbf{Shariah Status} \\
\midrule
CEX perpetual & Binance, Bybit, OKX & \mbox{\textcolor{ribared}{$0.01\%$/8h}} & No
  & \textcolor{ribared}{\textit{Riba}-bearing} \\[3pt]
Prem-only DEX & dYdX~v4, D8X & \mbox{\textcolor{halalgreen}{None}} & No
  & \textcolor{halalgreen}{Halal candidate} \\[3pt]
Hybrid DEX & Hyperliquid & \mbox{\textcolor{ribared}{$0.01\%$/8h}} & No
  & \textcolor{ribared}{\textit{Riba}-bearing} \\[3pt]
DEX expiring & MCDEX, Serum & \mbox{\textcolor{halalgreen}{None}} & Yes
  & \textcolor{halalgreen}{Strongest} \\
\bottomrule
\end{tabular}
\end{table}

Two entries deserve particular comment.
\textbf{Hyperliquid} is structurally a DEX --- non-custodial, on-chain
settlement --- yet its funding formula is identical to Binance's
\cite{Hyperliquid2025}. Its Shariah status is therefore identical to Binance's.
This exposes the literature's architectural error: calling something a DEX does
not eliminate \textit{riba}. Reading the formula does.

\textbf{DEX expiring futures} represent the theoretically strongest Shariah
case within the perpetual/futures universe, combining three favourable
properties simultaneously: (i)~no interest parameter ($I$ is structurally
inapplicable --- there is no periodic funding mechanism); (ii)~a fixed
settlement date, which directly addresses the \textit{gharar} objection about
perpetual indefiniteness that conservative scholars raise; and (iii)~smart
contract self-custody with no exchange counterparty risk. Where premium-only
DEX perpetuals satisfy the \textit{riba} criterion but require a rebuttal of
the perpetual \textit{gharar} objection, DEX expiring futures satisfy both
simultaneously. Protocols implementing or having implemented on-chain expiring
futures include MCDEX and Serum-based order books on Solana; liquidity
fragmentation across expiry dates remains the principal practical obstacle to
wider adoption. This paper focuses on the empirically active perpetuals market,
but Section~\ref{sec:limitations} situates DEX expiring futures as an important
direction for both market development and future Shariah research.

\subsection{Everlasting Options: A Forward-Looking Extension}
\label{sec:everlasting}

Ackerer, Hugonnier \& Jermann \cite{Ackerer2024} extend their perpetual futures
framework to derive pricing for \textbf{everlasting options} --- perpetual
payoff instruments that carry the right (but not the obligation) to exercise a
defined payoff at any point, with no fixed expiry date. Like perpetual futures,
everlasting options are sustained by a periodic premium paid by one party to the
other, eliminating the need for a settlement date while preserving a well-defined
payoff structure.

The premium-only pricing rule ($\iota = 0$) applies directly. Under the
Ackerer et al.\ framework, the no-arbitrage price of an everlasting option is:
\[
f_{o,t} = \mathrm{E}\!\left[\phi(x_{\theta_t})\right],
\]
where $\phi(\cdot)$ is the payoff function, $x_{\theta_t}$ is the underlying
asset price at the random stopping time $\theta_t$, and the expectation is taken
under the risk-neutral measure. When $\iota = 0$, no interest enters the pricing
equation: the premium paid between counterparties is purely market-driven,
determined by the option's time value and the underlying's dynamics rather than
any fixed financing charge.

From a Shariah perspective, everlasting options on RWA-backed asset tokens
possess several potentially favourable properties. First, the payoff is
\emph{defined} --- the option buyer's maximum loss is the premium paid, and the
payoff function $\phi$ is contractually specified. This bounded, asymmetric
risk profile addresses \textit{gharar} concerns more directly than perpetual
futures, where both gain and loss are unbounded. Second, when priced with
$\iota = 0$, no fixed predetermined charge exists, satisfying El-Gamal's
\textit{riba} conditions. Third, everlasting options on commodity tokens
(gold, silver, Shariah-compliant commodity indices) could enable halal hedging
with more defined risk profiles than perpetual futures --- the buyer knows at
entry the maximum loss, while the seller is compensated by market premium, not
interest.

This is explicitly a \textbf{forward-looking extension}. No \textit{fatwa} or
formal Shariah ruling currently addresses everlasting options. The standard
prohibitions on conventional options --- rooted in the uncertainty of the
premium and the sale of a non-existent asset --- may require re-examination when
applied to the everlasting structure, where the premium is market-determined and
no principal is transferred without delivery of the contractual payoff. Dedicated
jurisprudential research is warranted before any institutional implementation.
We flag this as the most theoretically significant near-term extension of the
halal perpetuals research agenda, alongside DEX expiring futures
(Section~\ref{sec:limitations}).

%% ============================================================
\section{Literature Review}
\label{sec:lit}
%% ============================================================

The 38 sources reviewed span five domains: (1) perpetual futures mechanics and
origin \cite{Hayes2016,Shiller1993,Ackerer2024}; (2) \textit{riba} in
derivatives \cite{ElGamal2006,Kamali2000,Ayub2007,Jobst2007,SoleJobst2012,
Malkawi2014,Tashfeen2022,Uluyol2024}; (3) \textit{gharar} in derivatives
\cite{ElGamal2001,Alias2015,Elhessi2018,Kasri2016,AliasOthman2024}; (4)
\textit{maysir} in derivatives \cite{Salamon2015,Alias2021,Alias2023,
Obaidullah2005}; and (5) contemporary \textit{fatwas} and institutional
standards \cite{Musaffa2024,IslamQA2023,SeekersGuidance2025,PracticalIF2023,
AAOIFI2015,DSNMUI2021,Amin2024,Indonesian2024}.
Across all 38 sources, zero papers identify, name, or compare dYdX~v4's
premium-only formula to the CEX interest-bearing formula.
This uniformity of the gap is itself the central finding of this review.

\subsection{The Progressive-Theoretical Tradition}

The progressive tradition in Islamic derivatives scholarship assembled, over
two decades, a set of analytical frameworks that collectively imply the
permissibility of premium-only DEX perpetuals --- without any author knowing
it, because the technology did not yet exist.

\textbf{Foundations of the riba analysis.}
El-Gamal \cite{ElGamal2006} provides the most rigorous economic definition of
\textit{riba}: a \emph{predetermined excess return that is independent of the
underlying economic outcome}.
Applied to the CEX funding formula, the interest floor $I = 0.01\%$ per 8 hours
is predetermined, fixed, and independent of whether the position profits ---
satisfying all three of El-Gamal's riba conditions.
Applied to dYdX~v4's premium-only formula ($I = 0$), none of the three
conditions is satisfied: the premium is market-determined, variable, and
can be negative.
El-Gamal's 2006 framework, rigorously applied, already yields our 2025 verdict ---
it simply had no instrument to apply it to at the time of writing.
Uddin \cite{Uddin2015} and Obaidullah \cite{Obaidullah2005} provide complementary
foundational definitions: \textit{riba} is any predetermined excess over principal,
regardless of label. The Binance API's own label of \texttt{interestRate} removes
any interpretive ambiguity.

\textbf{Frameworks for Shariah-compliant derivatives.}
Sole \& Jobst \cite{SoleJobst2012} develop a four-criteria framework for
Shariah-compliant derivatives, the first of which is an explicit zero-interest
parameter ($\iota = 0$).
Their 2012 paper articulates exactly the condition that dYdX~v4 satisfies in 2023
--- twelve years before the platform launched.
Jobst \cite{Jobst2007} independently proposed riba-free derivative structures based
on risk-sharing rather than fixed interest; our paper is the empirical
realisation of his theoretical blueprint.
Ayub \cite{Ayub2007} contributes a utility-based three-part permissibility test ---
no riba, no excessive gharar, legitimate economic function --- that DEX perpetuals
satisfy on all three counts: $I = 0\%$, on-chain transparency, and demonstrated
price-discovery function.

\textbf{The \textit{gharar} tradition.}
El-Gamal \cite{ElGamal2001} reframes \textit{gharar} not as all uncertainty but as
\emph{exploitative uncertainty rooted in information asymmetry}.
Under this definition, DEX protocols exhibit zero prohibited gharar: every
participant observes the same on-chain price feed, the same immutable funding
formula, and the same publicly auditable smart contract simultaneously.
Kamali \cite{Kamali2000} argues that standardisation and contractual transparency
are the decisive \textit{gharar} reducers in exchange-traded derivatives.
DEX smart contracts are more standardised --- and more certainly standardised ---
than any regulated exchange, because the rules are immutable code, not regulatory
policy subject to management discretion.
Elhessi \cite{Elhessi2018} demonstrates empirically that institutional transparency
reduces \textit{gharar} to permissible levels in Malaysian commodity futures;
blockchain immutability provides an analytically superior form of the same
transparency without reliance on institutional compliance.
Alias \& Othman \cite{AliasOthman2024}, in the most recent systematic gharar
review (2024), develop a five-dimension \textit{gharar} framework applicable to
crypto derivatives --- but apply it only to CEX platforms, reaching prohibition.
Applied to DEX perpetuals, every dimension of their framework yields a materially
different verdict: on-chain pricing eliminates information asymmetry,
automated liquidation removes counterparty uncertainty, and the public funding
formula removes settlement ambiguity.

\textbf{The \textit{maysir} tradition.}
Salamon \cite{Salamon2015} develops the most rigorous five-element maysir test:
(1) zero productive function, (2) pure chance, (3) one-sided gain, (4)
social harm, (5) wealth transfer without value creation.
DEX perpetuals used for hedging satisfy none of these criteria: they serve a
productive economic function, require knowledge and analysis, can benefit both
parties through risk transfer, improve market efficiency, and facilitate risk
redistribution rather than wealth destruction.
Alias \cite{Alias2021} applies a six-criteria maysir framework and reaches a
prohibitive conclusion --- but his criteria, applied to hedging contexts on
premium-only DEX platforms, do not support prohibition.

\textbf{The call for halal hedging instruments.}
Tashfeen \cite{Tashfeen2022} surveys the literature and explicitly calls for the
development of Shariah-compliant hedging instruments, citing riba elimination as
the key structural requirement.
This call was answered by technology within one year of publication: dYdX~v4
launched in 2023 with $I = 0\%$.
Amin \cite{Amin2024} and Uluyol \cite{Uluyol2024} take the final theoretical
step, explicitly arguing that derivatives without predetermined charges can be
halal --- without knowing that such instruments already existed.
Our paper is the empirical instantiation of every framework in this tradition:
we convert conditional theoretical possibilities into documented empirical facts.

\subsection{The Conservative-Prohibitive Tradition}

The conservative tradition is internally consistent and, within its domain of
application, correct.
Its error is categorical, not analytical: it treats the CEX perpetual as the
perpetual, and applies platform-specific findings to the instrument class.

\textbf{The foundational consensus.}
Obaidullah \cite{Obaidullah2005} provides the foundational conservative textbook
position: all interest-bearing components in derivatives are \textit{riba}; all
excessive uncertainty constitutes \textit{gharar}; instruments serving no productive
function are \textit{maysir}.
This framework is correct as applied to CEX perpetuals and our paper reaffirms it
for that platform type.
Our extension is narrow: the consensus was formed on instruments that universally
contained interest. We document a 2023 instrument that does not.

\textbf{The gharar analysis.}
Alias \cite{Alias2015} identifies information asymmetry and price uncertainty as the
operative \textit{gharar} in futures contracts and concludes they are prohibited.
On-chain smart contracts --- where pricing formulas are publicly auditable and
non-manipulable --- directly address Alias's operative concern.
Kasri \cite{Kasri2016} demonstrates that regulatory compliance frameworks
systematically fail to eliminate \textit{gharar} in practice, as structural gaming
persists.
His conclusion supports our argument: only structural design changes ($I = 0$ in
the funding formula) provide genuine Shariah compliance --- not trading rules or
usage policies imposed on top of riba-bearing CEX infrastructure.

\textbf{The contemporaneous failure.}
The most instructive cases are papers published in 2023 and 2024 --- after
dYdX~v4 launched --- that reach universal prohibitive conclusions with no
awareness of the new instrument.
Alias \cite{Alias2023} updates his prohibitive verdict in the same year dYdX~v4
launched, without mentioning it.
Musaffa \cite{Musaffa2024} writes in 2024 --- one year after dYdX~v4 --- and
correctly identifies the Binance interest floor as riba, then concludes all
perpetuals are prohibited.
IslamQA \cite{IslamQA2023} and Practical IF \cite{PracticalIF2023} commit the
same contemporaneous error.
Zulkarnain \& Syibly \cite{Indonesian2024} provide the most technically rigorous
conservative analysis, independently confirming that CEX funding ($F = P + \clamp(I-P)$,
$I = 0.01\%$) constitutes riba.
Their analysis is correct for its target --- Binance --- but is titled as analysis
of ``futures trading'' generally, the clearest instantiation of the categorical
extrapolation error.
These papers validate our CEX diagnosis and, in so doing, make the gap
explicit: every conservative analysis published through early 2025 treats the
Binance perpetual as the only perpetual, leaving the question of alternative
formula structures entirely unexamined.

\subsection{The Pivotal Concession: SeekersGuidance (2025)}

The most significant development in the conservative literature is
SeekersGuidance's 2025 ruling \cite{SeekersGuidance2025}. The ruling is
structured as a response to the specific question: \emph{``Are crypto perpetual
contracts still haram if there's no interest, and the asset price tracks?''}
The answer maintains prohibition --- but on narrowed grounds that are decisive
for our argument.
The ruling explicitly concedes that removing the interest component
\emph{resolves the riba objection}, then identifies the remaining prohibition
as principally \textit{qabdh} --- the classical possession requirement under
traditional fiqh.
This concession is jurisprudentially significant for two reasons.
First, the most rigorous contemporary conservative authority now agrees that
\textit{riba}-free perpetuals are not prohibited by \textit{riba} --- validating
our central structural claim.
Second, the remaining objection (\textit{qabdh}) is narrower and more tractable
than riba: it is addressed by real-world asset (RWA)-backed tokens where physical
custody is established, by the \textit{istisna}/\textit{salam} precedent of
deferred delivery, and by DSN-MUI Fatwa 144 \cite{DSNMUI2021} classifying
crypto assets as Shariah-legitimate commodities.
The conservative literature has, in its most recent development, converged on
precisely the question our paper answers.

\subsection{The Gap: Uniform and Structural}

The gap in the 38-source literature is structurally uniform: no paper ---
progressive or conservative, pre-2020 or published in 2024 --- identifies,
names, or compares the CEX interest-bearing formula to any alternative.
The progressive tradition built frameworks that logically imply DEX permissibility;
the conservative tradition reached verdicts that correctly apply to CEX and are
now explicitly validated by SeekersGuidance for that platform.
Both traditions analysed the same instrument --- the Binance-style perpetual ---
and reached opposite conclusions about its permissibility.
Neither noticed that a structurally different instrument, with $I = 0$
documented in its official protocol specification, had emerged and was operating
at scale.

The most recent DEX-focused Islamic finance scholarship \cite{ResearchGate2026}
analyses spot trading on decentralised exchanges and concludes that DEX spot
transactions can satisfy Islamic sale-and-purchase principles --- but does not
extend the analysis to perpetual derivatives.
Our paper makes the necessary extension: the same on-chain transparency and
non-custodial structure that makes DEX spot trading potentially halal is equally
operative in DEX perpetuals, with the additional decisive feature that premium-only
perpetuals eliminate the interest parameter that makes spot-equivalent CEX
derivatives haram.

A wave of proactive-framing scholarship published between 2024 and 2026 validates
the methodological approach of this paper. Ghazali et al.\ \cite{Ghazali2025}
explicitly advocate for ``structured alternatives replicating functions while
adhering to Shariah'' --- precisely the design logic of premium-only DEX
perpetuals. Alsmadi et al.\ \cite{Alsmadi2025} call for ``a proactive
multi-faceted approach to align digital innovation with Shariah,'' documenting
how Islamic Fintech firms that design compliance into their architecture
outperform those that apply compliance as a post-hoc filter. Farhat
\cite{Farhat2025} demonstrates that ``Shariah compliance through deliberate
screening and structuring'' --- as opposed to avoidance --- systematically
improves firm performance and risk profiles, suggesting that structured
participation in perpetual markets could benefit Muslim investors. Modjo et al.\
\cite{Modjo2025} show that ``digital innovation amplifies the benefits of
Shariah-compliant practices,'' providing empirical evidence that technology-native
Shariah compliance (exactly what premium-only DEX protocols represent) generates
superior outcomes compared to non-digital compliant structures.
S\&P Global \cite{SP2025} documents a \$190--200 billion sukuk market built
entirely on premium and profit-sharing structures without interest --- the
largest empirical proof-of-concept that market-scale Islamic finance is achievable
without predetermined interest charges. Together, these five sources confirm that
the field has moved decisively toward proactive design framing: the question is
no longer whether interest can be structurally eliminated, but how to apply that
elimination to new instrument classes. This paper applies it to perpetual futures
for the first time.

The contribution of this paper is to make the connection that 38 prior sources,
including the most recent and most technically sophisticated, did not make:
the CEX/DEX formula distinction is a Shariah distinction, the progressive
frameworks already imply DEX permissibility, and the conservative prohibitions
remain valid --- but only for CEX perpetuals.

%% ============================================================
\section{Empirical Evidence: A Natural Experiment}
\label{sec:empirical}
%% ============================================================

\subsection{Design and Data}

We collected 365-day funding rate data (February 2025 -- February 2026) from
five platforms, yielding 40,218 intervals (Table~\ref{tab:stats}). The dataset
constitutes a \textbf{natural experiment}: within the DEX category, two
platforms implement different values of $I$ --- dYdX~v4 with $I=0$ and
Hyperliquid with $I=0.01\%$/8h. Both are DEXs, operating in the same market,
over the same period. The only controlled variable is the interest parameter.
If the interest parameter matters, the two DEXs should produce different
distributional outcomes. If architecture matters instead, they should produce
similar outcomes.

All data were collected via public APIs at each platform's native interval:
8-hourly for CEX; hourly for dYdX~v4 and Hyperliquid (Hyperliquid's 8-hour
equivalent intervals are compared by summing three consecutive hourly payments).
Figure~\ref{fig:timeseries} shows the complete 365-day time series.

\begin{table*}[htbp]
\centering
\caption{Funding rate descriptive statistics, 365-day dataset (February 2025 --
February 2026). CEX and Hyperliquid figures are per 8-hour interval. dYdX~v4
figures are per native 1-hour interval. Shariah classification is determined by
the interest parameter $I$, not by CEX/DEX architecture.}
\label{tab:stats}
\small
\setlength{\tabcolsep}{5pt}
\begin{tabular}{@{}llcrrrrcc@{}}
\toprule
\textbf{Platform} & \textbf{Asset} & \textbf{Category} & \textbf{$N$} &
\textbf{Mean} & \textbf{Std} & \textbf{Min} & \textbf{Max} & \textbf{Neg.\%} \\
\midrule
Binance & BTC & CEX & 1{,}095 & 0.0041\%/8h & 0.0043\% & $-$0.015\% & $+$0.010\% & 15.6\% \\
Binance & ETH & CEX & 1{,}095 & 0.0037\%/8h & 0.0050\% & $-$0.037\% & $+$0.010\% & 19.7\% \\
Bybit   & BTC & CEX & 1{,}094 & 0.0040\%/8h & 0.0050\% & $-$0.029\% & $+$0.011\% & 21.7\% \\
Bybit   & ETH & CEX & 1{,}094 & 0.0037\%/8h & 0.0067\% & $-$0.118\%$^{a}$ & $+$0.010\% & 25.9\% \\
OKX     & BTC & CEX &   318   & 0.0032\%/8h & 0.0040\% & $-$0.009\% & $+$0.010\% & 20.1\% \\
OKX     & ETH & CEX &   318   & 0.0007\%/8h & 0.0050\% & $-$0.022\% & $+$0.010\% & 40.9\% \\
\addlinespace
\textit{CEX pooled} & & & \textit{5{,}014} & \textit{0.0037\%/8h} & & & & \textit{22.0\%} \\
\midrule
dYdX~v4 & BTC & Prem-DEX ($I=0$) & 8{,}842 & 0.0006\%/1h & 0.0019\% & $-$0.013\% & $+$0.054\% & 30.8\% \\
dYdX~v4 & ETH & Prem-DEX ($I=0$) & 8{,}842 & 0.0010\%/1h & 0.0019\% & $-$0.014\% & $+$0.022\% & 20.6\% \\
\addlinespace
\textit{dYdX pooled} & & & \textit{17{,}684} & \textit{0.0008\%/1h} & & & & \textit{25.7\%} \\
\midrule
Hyperliquid & BTC & Hybrid-DEX ($I>0$) & 8{,}760 & 0.0011\%/1h & 0.0013\% & $-$0.009\% & $+$0.023\% & 11.8\% \\
Hyperliquid & ETH & Hybrid-DEX ($I>0$) & 8{,}760 & 0.0009\%/1h & 0.0015\% & $-$0.039\% & $+$0.015\% & 18.2\% \\
\addlinespace
\textit{Hyp. pooled} & & & \textit{17{,}520} & \textit{0.0010\%/1h} & & & & \textit{15.0\%} \\
\bottomrule
\multicolumn{9}{l}{\footnotesize $^{a}$Single-event extreme; acute market backwardation episode. $N$ (total) $= 40{,}218$.}
\end{tabular}
\end{table*}

\begin{figure*}[htbp]
\centering
\includegraphics[width=\textwidth]{fig_funding_timeseries.png}
\caption{365-day funding rate time series across all five platforms (February
2025 -- February 2026), BTC (top) and ETH (bottom). The dashed line at
$+0.01\%$ marks the CEX interest floor. CEX rates (Binance, Bybit, OKX) cluster
persistently above zero, anchored by $I$. dYdX~v4 oscillates symmetrically
around zero, consistent with $I=0$. Hyperliquid --- plotted here at its native
hourly rate, $8\times$ smaller than the 8h CEX intervals shown --- exhibits a
positive bias consistent with its $I=0.01\%$/8h formula, mirroring CEX
behaviour rather than dYdX.}
\label{fig:timeseries}
\end{figure*}

\subsection{The Natural Experiment: Results}

\subsubsection{The Primary Test: CEX vs.\ Premium-only DEX}

Table~\ref{tab:tests} presents statistical tests comparing CEX (Binance + Bybit
+ OKX) against dYdX~v4. The null hypothesis is that the two distributions are
drawn from the same population. All three tests reject decisively.

\begin{table}[htbp]
\centering
\caption{Statistical tests: CEX ($n = 5{,}014$, per 8h) vs.\ dYdX~v4
($n = 17{,}684$, per hour). Hyperliquid is excluded from this comparison as it
uses $I = 0.01\%$/8h.}
\label{tab:tests}
\small
\begin{tabular}{lrc}
\toprule
\textbf{Test} & \textbf{Statistic} & \textbf{$p$-value} \\
\midrule
Two-sample $t$-test     & $t = 59.95$     & $< 10^{-300}$ \\
Mann-Whitney $U$        & $U = 64{,}005{,}833$ & $< 10^{-300}$ \\
Kolmogorov-Smirnov      & $D = 0.487$     & $< 10^{-300}$ \\
\midrule
Cohen's $d$ (effect size) & $d = 0.725$ & Large \\
\bottomrule
\end{tabular}
\end{table}

The KS statistic of 0.487 means that nearly half the mass of the two
distributions does not overlap --- a structural divergence, not a mean shift.
Cohen's $d = 0.725$ is a large effect size by any standard, and it persists
across the full 365-day window covering both bull and bear regimes.

\subsubsection{The Control: Hyperliquid vs.\ dYdX}

The natural experiment's control condition answers the architectural question
directly. If DEX structure were the operative variable, Hyperliquid and dYdX~v4
--- both DEXs, same market, same period --- should produce similar distributions.
They do not.

\begin{itemize}[noitemsep, leftmargin=*]
\item \textbf{Negative-rate frequency:} dYdX~v4 25.7\% vs.\ Hyperliquid 15.0\%.
The platform with $I=0$ shows \emph{higher} negative frequency, consistent with
a symmetric mechanism unanchored from zero. Hyperliquid, with $I > 0$, shows
lower negative frequency --- the interest floor creates an upward bias identical
in character to CEX.
\item \textbf{Distribution shape:} Hyperliquid's distribution clusters with CEX
(both positively biased), not with dYdX~v4 (symmetric). Architecture does not
determine shape; the interest parameter does.
\end{itemize}

This is the most powerful result in the paper. Two DEXs. Same market. Same
year. Different formulas. Different distributions. The difference is the
interest parameter.

\subsubsection{Distributional Evidence}

Figure~\ref{fig:dist} shows the distribution comparison across the three
empirically observed perpetuals categories (DEX expiring futures are excluded
as no 365-day dataset was available for this study).
Figure~\ref{fig:boxplot} shows per-platform distributions. Figure~\ref{fig:neg}
shows negative-rate frequencies, which provide the clearest visual signature of
the interest floor: platforms with $I > 0$ are suppressed in the negative tail
precisely because the interest component offsets downward premiums.

\begin{figure}[htbp]
\centering
\includegraphics[width=\columnwidth]{fig_funding_distribution.png}
\caption{Funding rate distributions by category. CEX (red, $n=5{,}014$, per 8h
interval) is skewed positive by the interest floor. dYdX~v4 (green,
$n=17{,}684$, per hour) is centred near zero, consistent with $I=0$. Hyperliquid
(amber, $n=17{,}520$) clusters with CEX despite its DEX architecture. Cohen's
$d = 0.725$, KS $= 0.487$ (CEX vs.\ dYdX).}
\label{fig:dist}
\end{figure}

\begin{figure}[htbp]
\centering
\includegraphics[width=\columnwidth]{fig_platform_boxplot.png}
\caption{Platform-level funding rate distributions. CEX platforms (Binance,
Bybit, OKX) and Hyperliquid share a positive-biased structure consistent with
$I > 0$. dYdX~v4 is more symmetric, consistent with $I = 0$.}
\label{fig:boxplot}
\end{figure}

\begin{figure}[htbp]
\centering
\includegraphics[width=\columnwidth]{fig_negative_funding.png}
\caption{Negative funding rate frequency by platform (BTC, 365 days).
dYdX~v4 ($\approx$25.7\%, $I=0$) shows the highest negative-rate frequency
--- the signature of an unconstrained symmetric mechanism. Hyperliquid
(11.8\%, $I>0$) shows the lowest --- the interest floor suppresses negative
rates below the CEX average (22.0\% pooled), confirming the formula dominates
over architecture.}
\label{fig:neg}
\end{figure}

\subsection{The Structural vs.\ Observed Cost Distinction}

A crucial Shariah distinction operates at the cost level: the \emph{structural
riba} component versus the \emph{observed market premium} component.

\begin{itemize}[noitemsep, leftmargin=*]
\item \textbf{CEX structural riba} (unavoidable, regardless of market):
$0.01\% \times 3 \times 365 = \mathbf{10.95\%}$/year. Even in a falling market
with negative premiums, this floor continues.

\item \textbf{CEX observed total} (interest + market premium, our 365-day data):
$0.0037\% \times 3 \times 365 \approx \mathbf{4.0\%}$/year. Below 10.95\%
because negative premiums partially offset interest during bear phases --- but
the \textit{riba} component is present in every interval.

\item \textbf{dYdX structural riba}: $\mathbf{0\%}$/year. No interest
component in the formula.

\item \textbf{dYdX observed total} (market premium only, our 365-day data):
$0.0008\% \times 24 \times 365 \approx \mathbf{6.8\%}$/year during the 2025
bull period. This number exceeded CEX's observed total because positive premiums
were large --- but every dollar of this premium is bilateral market transfer
between traders, not interest. The Shariah position of a Muslim paying 6.8\%
in halal premium is categorically different from one paying 4.0\% in which an
inextricable riba component is embedded.
\end{itemize}

Figure~\ref{fig:cost} illustrates both components. The Shariah verdict does not
depend on which number is larger: it depends on whether a fixed predetermined
charge is present at all. It is; and it is not.

\begin{figure}[htbp]
\centering
\includegraphics[width=\columnwidth]{fig_cost_comparison.png}
\caption{Left: structural interest burden (riba component only). CEX and
Hyperliquid impose 10.95\%/year regardless of market; dYdX~v4 imposes 0\%.
Right: observed total funding cost over the 365-day period. dYdX~v4 exceeded
CEX in this bull period, but the dYdX cost is halal market premium, not riba.
Shariah compliance is determined by the left panel, not the right.}
\label{fig:cost}
\end{figure}

%% ============================================================
\section{Shariah Analysis}
\label{sec:shariah}
%% ============================================================

We apply the three canonical Islamic finance prohibitions across all four
platform categories. Table~\ref{tab:verdicts} summarises the verdicts.
Primary analysis focuses on Premium-only DEX perpetuals (dYdX~v4 type, $I=0$)
as the empirically validated category; DEX expiring futures are discussed
as a theoretically superior alternative pending liquidity development.

\begin{table}[htbp]
\centering
\caption{Shariah verdicts by category and prohibition.}
\label{tab:verdicts}
\footnotesize
\setlength{\tabcolsep}{3pt}
\begin{tabular}{@{}p{1.55cm}cccc@{}}
\toprule
\textbf{Category} & \textbf{Riba} & \textbf{Gharar} & \textbf{Maysir} & \textbf{Overall} \\
\midrule
CEX perpetual & \textcolor{ribared}{Haram} & Concern & Cond. & \textcolor{ribared}{Haram} \\
Hybrid DEX & \textcolor{ribared}{Haram} & Reduced & Cond. & \textcolor{ribared}{Haram} \\
Prem-only DEX & \textcolor{halalgreen}{Clear} & Reduced & Cond. & \textcolor{halalgreen}{Candidate} \\
DEX expiring & \textcolor{halalgreen}{N/A} & \textcolor{halalgreen}{Minimal} & Cond. & \textcolor{halalgreen}{Strongest} \\
\bottomrule
\multicolumn{5}{@{}l@{}}{\scriptsize Cond.\ = context-dependent on trader intent and leverage.}\\
\multicolumn{5}{@{}l@{}}{\scriptsize DEX expiring futures: no funding mechanism; fixed settlement date.}\\
\end{tabular}
\end{table}

\subsection{Riba: The Decisive Prohibition}

\textit{Riba} is the decisive prohibition because it operates structurally,
not conditionally. El-Gamal \cite{ElGamal2006} operationalises \textit{riba}
through three conditions: (i)~the excess is predetermined; (ii)~it is
independent of the underlying economic outcome; (iii)~it is fixed in advance.
CEX interest ($I = 0.01\%$/8h) satisfies all three. dYdX~v4 premium satisfies
none:

\begin{itemize}[noitemsep, leftmargin=*]
\item The premium is determined by the market at the time of payment --- not
predetermined.
\item It depends entirely on the relative price of perpetual versus spot --- not
independent of economic outcome.
\item It can be zero, positive, or negative --- not fixed in advance.
\end{itemize}

Sole \& Jobst \cite{SoleJobst2012} specify the $\iota=0$ interest criterion as
the primary condition for Shariah-compliant derivatives. dYdX~v4 satisfies it.
SeekersGuidance (2025) \cite{SeekersGuidance2025} concedes that satisfying this
criterion resolves the \textit{riba} objection. The \textit{riba} verdict for
premium-only DEX perpetuals is clear.

\subsection{Gharar: Substantially Reduced}

The textual foundation of the \textit{gharar} prohibition is the well-known
hadith recorded in \textit{Sahih Muslim}: ``The Prophet (peace be upon him) forbade
the sale of \textit{gharar}'' (\textit{Sahih Muslim}, Book~10, No.~3614).
Classical jurists operationalised this prohibition as covering contracts with
major ambiguity in terms, price, or delivery that creates the potential for
exploitation or dispute. \textit{Gharar} --- prohibited uncertainty arising
from informational asymmetry or ambiguous contract terms --- is therefore a
legitimate concern in derivatives \cite{ElGamal2001,Alias2015,AliasOthman2024}.
The question is not whether any uncertainty exists (all commerce involves
uncertainty) but whether the uncertainty is of the prohibited exploitative type
\cite{ElGamal2001}.

A specific subcategory relevant to perpetuals is \textit{bay' ma'dum} --- the
sale of a non-existent or undeliverable asset (Hadith, \textit{Sunan
al-Tirmidhi}, No.~1232: ``Do not sell what you do not have'').
Cash-settled CEX perpetuals raise this concern acutely: the trader never
possesses the underlying asset, and settlement is in a numeraire disconnected
from physical delivery.
DEX perpetuals written on RWA-backed tokens, where the token represents a
verified legal claim on a physical asset held in custody, directly address the
\textit{bay' ma'dum} concern by tying the position to a genuinely existent,
deliverable underlying.

DEX smart contracts address the broader \textit{gharar} concerns more
completely than any regulated exchange:

\begin{itemize}[noitemsep, leftmargin=*]
\item \textbf{Contractual terms} are published in immutable on-chain code,
publicly auditable by any party. No information asymmetry exists between
participants and protocol.
\item \textbf{Price discovery} occurs through public order books or AMM
formulas with no privileged access.
\item \textbf{Loss limitation} is automatic: positions are liquidated at a
specified margin threshold. A trader cannot lose more than deposited collateral.
This directly addresses Obaidullah's \cite{Obaidullah2005} concern about
unlimited potential loss and eliminates the unbounded-loss \textit{gharar}
that characterises uncollateralised OTC contracts.
\end{itemize}

El-Gamal \cite{ElGamal2001} defines prohibited \textit{gharar} as exploitative
uncertainty: one party exploiting superior information against another. DEX
protocols achieve zero information asymmetry by design --- every participant
observes the same on-chain price feed, the same funding formula, and the same
contract terms simultaneously. Kamali \cite{Kamali2000} holds that
standardisation resolves \textit{gharar}; DEX smart contract rules are more
standardised and more immutable than any regulatory standard, because they are
mathematically enforced rather than institutionally promised.

We conclude that \textit{gharar} is substantially reduced in DEX perpetuals
relative to conventional futures, falling below the threshold of prohibition
under the El-Gamal and Kamali frameworks. The residual \textit{gharar} concern
--- oracle risk in price feed computation --- is noted in Section~\ref{sec:limitations}
and warrants separate Shariah analysis as the oracle ecosystem matures.

\subsection{Maysir: Conditional on Use}

The \textit{maysir} prohibition is grounded in an explicit Quranic injunction:
``O believers! Intoxicants, gambling (\textit{maysir}), idols, and drawing lots
are all evil of Satan's handiwork --- so shun them so you may be successful''
(Qur'an 5:90). Classical jurists interpreted \textit{maysir} as any
transaction in which wealth transfers through pure chance without productive
economic function --- the paradigm case being gambling where one party gains
only at another's direct loss, with no underlying value created.

\textit{Maysir} in a derivatives context requires zero productive economic
function, pure chance determination, and one-sided wealth transfer
\cite{Salamon2015,Alias2021}. This prohibition is use-contingent rather than
structural: the same instrument can constitute \textit{maysir} in a speculative
context and legitimate commerce in a hedging context.

DEX perpetuals used for hedging fail the \textit{maysir} criteria on all
counts: (1)~price discovery and risk transfer serve productive economic
functions \cite{Shiller1993}; (2)~outcomes depend substantially on skill,
analysis, and information --- not pure chance; (3)~the transfer is bilateral
--- shorts benefit when longs pay and longs benefit when shorts pay, depending
on market direction. The zero-sum critique raised by Practical IF
\cite{PracticalIF2023} applies equally to insurance contracts and secondary
market equity trading, neither of which classical scholars prohibit as
\textit{maysir}.

Critically, intent (\textit{niyyah}) matters. Pure speculation on short-term
price movements without economic basis more closely approximates \textit{maysir}
than disciplined hedging with documented rationale. The instrument does not
compel either use; the trader's intent and discipline determine the verdict.

%% ============================================================
\section{Conservative Objections and Rebuttals}
\label{sec:objections}
%% ============================================================

We address five objections systematically. Where possible, we cite the
objection's own scholarship to demonstrate that the concern is addressed rather
than dismissed.

\subsection{Objection 1: Riba Persists Through Leverage and Implicit Credit}

\emph{``Even without the explicit interest parameter, leverage implies a loan
that bears implicit interest.''} --- SeekersGuidance \cite{SeekersGuidance2025};
IslamQA \cite{IslamQA2023}.

\textbf{Response.} In DEX perpetuals, no loan is extended. The trader deposits
collateral (margin) and establishes a synthetic position. The protocol does not
advance credit; the collateral remains the trader's own funds posted against
a position. This is structurally equivalent to the classical \textit{hamish
jiddiyyah} (security deposit), not a loan. The premium payment, when it occurs,
is a bilateral market transfer between counterparties --- not a financing charge
from a lender \cite{SoleJobst2012,Malkawi2014}. SeekersGuidance (2025) itself
concedes this point in acknowledging that the riba objection is resolved when
no interest component is present.

\subsection{Objection 2: Gharar From Perpetual Nature and Price Uncertainty}

\emph{``The absence of a settlement date creates excessive gharar.''} ---
Musaffa \cite{Musaffa2024}; Alias \cite{Alias2015}.

\textbf{Response.} Perpetual nature does not constitute prohibited
\textit{gharar}. The contract terms at any given moment are fully specified: the
funding formula is hardcoded, the margin requirements are known, the liquidation
threshold is public. The absence of a settlement date is a feature, not an
ambiguity --- and indeed, the absence of an expiry removes the source of
\textit{gharar} that classical scholars identified in \textit{bai al-kali bil
kali} (deferred payment for deferred delivery). Automated liquidation bounds
loss to deposited margin; the trader cannot lose more than is risked
\cite{Kamali2000,Elhessi2018}.

\subsection{Objection 3: Maysir Through Zero-Sum Structure}

\emph{``One party wins only because another loses --- this is gambling.''}
--- Practical IF \cite{PracticalIF2023}; Alias \cite{Alias2021}.

\textbf{Response.} Every secondary market transaction has a buyer and a seller
holding opposing views; every insurance contract has a premium payer and a
beneficiary. If zero-sum payoff structure constitutes \textit{maysir}, then
all secondary market activity is prohibited --- a conclusion no classical scholar
has reached. The operative test is \emph{economic utility}, not aggregate payoff
symmetry \cite{Ayub2007,ElGamal2006}. Derivatives markets enable hedgers to
transfer unwanted risk to willing risk-takers, improving allocative efficiency.
Kamali \cite{Kamali2000} and Ayub \cite{Ayub2007} both permit instruments
serving this function.

\subsection{Objection 4: Non-Ownership and Qabdh (Possession) Requirement}

\emph{``Do not sell what you do not possess.''} --- IslamQA \cite{IslamQA2023};
SeekersGuidance \cite{SeekersGuidance2025}.

\textbf{Response.} This is the objection that SeekersGuidance (2025) identifies
as the \emph{remaining} concern after conceding the \textit{riba} point. Three
responses:

(a)~\textbf{RWA-backed tokens}: DEX perpetuals written on tokenised real-world
assets --- gold tokens, commodity-backed tokens --- where the token represents a
genuine legal claim on a physical asset held in custody satisfy the ownership
requirement directly. The DSN-MUI \cite{DSNMUI2021} recognition of
cryptocurrency as a commodity provides the institutional foundation.

(b)~\textbf{Historical precedent}: Classical scholars permitted \textit{salam}
and \textit{istisna}, both of which involve payment before delivery (deferred
\textit{qabdh}). The \textit{qabdh} requirement is not absolute; it yields to
established commercial necessity and contractual specificity.

(c)~\textbf{Evolving scholarly consensus}: The AAOIFI recognition of digital
assets \cite{AAOIFI2015} and the growing fiqh literature on electronic
\textit{qabdh} (constructive possession through exclusive, transferable digital
property rights) increasingly supports treating DEX token positions as
satisfying ownership requirements.

\subsection{Objection 5: Societal Harm Through Excessive Speculation}

\emph{``High-leverage derivatives cause financial instability and individual
ruin.''} --- DSN-MUI \cite{DSNMUI2021}; Kasri \cite{Kasri2016}.

\textbf{Response.} This objection applies to \emph{irresponsible use}, not to
the instrument itself. The same critique applies to margin lending in equities,
foreign exchange trading, and leveraged property investment --- none of which
scholars have categorically prohibited. Islamic scholars have long permitted
long-distance trade carrying risk of total loss when the risk serves productive
purposes. DEX perpetuals used with disciplined position sizing ($\leq 3\times$
leverage), algorithmic stop-losses, and asset selection focused on productive
underlying assets are no more harmful than established leveraged instruments.
The \textit{maqasid al-Shariah} objective of protecting wealth (\textit{hifz
al-mal}) is served by appropriate risk management, not by categorical exclusion.

Table~\ref{tab:objections} consolidates the objection-rebuttal structure.

\begin{table*}[htbp]
\centering
\caption{Conservative objections and DEX rebuttals.}
\label{tab:objections}
\small
\begin{tabular}{p{3.0cm}p{3.2cm}p{6.0cm}}
\toprule
\textbf{Objection} & \textbf{Source} & \textbf{Rebuttal} \\
\midrule
Riba from leverage / implicit credit &
SeekersGuidance (2025); IslamQA (2023) &
No credit extended; collateral $\neq$ loan; \textit{hamish jiddiyyah}
analogy; SeekersGuidance itself concedes riba is absent when $I=0$ \\[4pt]
\textit{Gharar} from no expiry and loss uncertainty &
Musaffa (2024); Alias (2015) &
Terms fully specified at all times; auto-liquidation caps loss to posted
margin; no deferred-delivery ambiguity \\[4pt]
\textit{Maysir} through zero-sum character &
Practical IF (2023); Alias (2021) &
Insurance and secondary equity markets are also zero-sum; operative test is
economic utility; hedging is not gambling \\[4pt]
Non-ownership and \textit{qabdh} &
SeekersGuidance (2025); IslamQA (2023) &
RWA-backed tokens satisfy \textit{qabdh}; \textit{salam}/\textit{istisna}
precedent; DSN-MUI commodity recognition \\[4pt]
Societal harm from leverage &
DSN-MUI (2021); Kasri (2016) &
Applies to irresponsible use; same critique fails for margin equity;
\textit{maqasid} served by prudent risk tools \\
\bottomrule
\end{tabular}
\end{table*}

%% ============================================================
\section{Implementation Framework for Muslim Traders}
\label{sec:implementation}
%% ============================================================

\subsection{Platform Selection: The Formula Test}

The single most important selection criterion is the interest parameter $I$ in
the funding formula, documented in each platform's official API and technical
documentation. Do not rely on marketing descriptions such as ``decentralised''
or ``non-custodial.''

\textbf{Currently compliant:} dYdX~v4 (\texttt{I=0}, documented \cite{dYdX2024})
and D8X (\texttt{I=0}, documented \cite{D8X2024}).

\textbf{Currently non-compliant despite DEX architecture:} Hyperliquid
(\texttt{I=0.01\%/8h}, documented \cite{Hyperliquid2025}). Using Hyperliquid
subjects a Muslim trader to the identical \textit{riba} burden as using Binance.

As new DEX platforms emerge, verify the funding formula documentation before
participation. Smart contract governance may alter parameters over time;
continuous monitoring is required.

\subsection{Asset Selection: Productive Underlying}

Prioritise perpetuals on assets with demonstrated economic utility:
\begin{itemize}[noitemsep, leftmargin=*]
\item \textbf{Commodity-backed tokens} (gold, silver, oil tokens) --- directly
relevant to \textit{qabdh} and \textit{maqasid} arguments.
\item \textbf{Infrastructure tokens} with significant productive usage (network
fee revenue, validator staking).
\item \textbf{Tokenised real-world assets} (RWA tokens representing verified
physical assets).
\end{itemize}
Avoid perpetuals on purely speculative or utility-less tokens. Asset selection
strengthens the maysir analysis by establishing genuine economic purpose
\cite{Kalimullina2020}.

\noindent\textit{Note on everlasting options:} For certain hedging applications
--- particularly those requiring a defined maximum loss rather than an uncapped
directional position --- everlasting options on RWA-backed commodity tokens
(Section~\ref{sec:everlasting}) may offer superior Shariah defensibility compared
to perpetual futures. The bounded payoff structure addresses \textit{gharar}
concerns more directly, and pricing without interest ($\iota = 0$) preserves
the structural \textit{riba}-free character. Pending dedicated scholarly review
and formal \textit{fatwa}, institutions with more conservative Shariah
requirements should note this instrument class as a potentially stronger
candidate for board approval than outright perpetual futures positions.

\subsection{Risk Management: Parameters}

Implement disciplined position management consistent with \textit{hifz al-mal}:
\begin{itemize}[noitemsep, leftmargin=*]
\item Maximum leverage: $2\text{--}3\times$ (not the $50\text{--}100\times$
offered by most platforms).
\item Mandatory stop-loss: maximum 10--15\% drawdown from entry before
automatic exit.
\item Position sizing: maximum 5\% of portfolio in any single perpetual position.
\item Regular rebalancing: review position rationale monthly.
\end{itemize}

\subsection{Intent Documentation (\textit{Niyyah})}

Document trading intent before opening positions: the economic rationale
(hedging, market-making, investment), the target asset, the risk parameters,
and the exit criteria. This serves both Shariah discipline and practical risk
management. A position opened without articulated rationale is more difficult
to distinguish from speculative gambling (\textit{maysir}) in retrospect.

\subsection{Scholarly Consultation}

This paper provides the structural analysis; it does not constitute a
\textit{fatwa}. Muslim traders and institutions should engage qualified Shariah
scholars in their jurisdiction with this technical analysis as input. We
particularly encourage engagement with AAOIFI and OIC Fiqh Academy on formal
standard-setting for DEX derivative instruments.

%% ============================================================
\section{Limitations and Future Research}
\label{sec:limitations}
%% ============================================================

\textbf{Single-year window.} Our 365-day dataset spans a full market cycle and
shows structural persistence, but multi-year data across multiple cycles would
further strengthen the empirical foundation. Future work should extend to
additional DEX platforms (D8X, Vertex) as their data matures.

\textbf{Platform governance risk.} DEX smart contracts are governed by token
holder votes. A governance proposal could introduce a non-zero $I$ parameter on
currently compliant platforms. This represents a monitoring requirement, not a
permanent compliance guarantee. We recommend institutional participants track
governance forums actively.

\textbf{Oracle risk.} The price feeds that compute the premium index on DEX
perpetuals introduce a residual \textit{gharar} concern: oracle manipulation
or failure could distort the funding rate. This risk is distinct from the
funding formula and warrants separate Shariah analysis as the oracle ecosystem
matures.

\textbf{Fatwa gap.} No recognised Shariah authority has issued a formal
\textit{fatwa} specifically addressing DEX perpetuals. The SeekersGuidance
(2025) ruling is the closest, and its concession on \textit{riba} is
significant, but formal institutional guidance would substantially advance
the field. This paper provides the technical foundation; the jurisprudential
formalisation must follow.

\textbf{DEX expiring futures.}
This paper's empirical analysis focuses on perpetual contracts because they
constitute the dominant liquid market. However, our taxonomy (Table~\ref{tab:taxonomy})
identifies DEX expiring futures as the theoretically superior Shariah category:
they combine zero interest (no funding mechanism exists at all), a fixed
settlement date (directly resolving the perpetual \textit{gharar} objection),
and smart-contract self-custody (no exchange counterparty risk). On-chain
expiring futures have been implemented on protocols including MCDEX and
Serum-based order books; their primary limitation is liquidity fragmentation
across expiry dates rather than any structural Shariah concern.
As on-chain liquidity infrastructure matures --- particularly through unified
expiry-date pools and cross-margining mechanisms --- DEX expiring futures may
offer the cleanest path to Shariah-compliant leveraged derivatives that also
satisfies the most conservative scholarly positions on \textit{gharar}.
Dedicated empirical and jurisprudential analysis of these instruments is
the most important near-term extension of this research agenda.

%% ============================================================
\section{Conclusion}
\label{sec:conclusion}
%% ============================================================

The Islamic finance consensus that perpetual futures are prohibited is correct
--- for CEX perpetuals. The error committed by 38 reviewed sources is an error
of extrapolation: correct diagnosis of \textit{riba} in one implementation
was incorrectly generalised to all implementations. That generalisation was
defensible as long as all implementations shared the same formula. It is no
longer defensible.

This paper establishes four findings that together constitute a new paradigm.

\textbf{First}, the interest parameter $I$ in CEX funding formulas is
demonstrable \textit{riba} by any classical definition --- predetermined, fixed,
labeled as interest, and unavoidable. Conservative scholars reached this verdict
correctly, including SeekersGuidance (2025), which provides the most rigorous
contemporary ruling. That verdict stands for all platforms with $I > 0$,
whether CEX (Binance, Bybit, OKX) or DEX (Hyperliquid).

\textbf{Second}, the interest parameter is mathematically unnecessary. Ackerer,
Hugonnier \& Jermann (2024) proved this formally. DEX protocols dYdX~v4 and
D8X proved it practically --- they have operated successfully with $I = 0$ at
scale. Riba was chosen in 2016; it was never required.

\textbf{Third}, 365 days and 40,218 funding intervals across a natural
experiment confirm that the interest parameter --- not exchange architecture ---
determines distributional structure. Hyperliquid (DEX, $I > 0$) behaves like
CEX. dYdX~v4 (DEX, $I = 0$) behaves differently. The Shariah significance of
this result is not statistical: it is the empirical proof that removing the
interest parameter changes the observable character of the instrument in exactly
the way Shariah analysis predicts.

\textbf{Fourth}, principles-based Shariah analysis --- applying El-Gamal's riba
conditions, Kamali's standardisation criterion, Sole \& Jobst's four-part
framework, and Salamon's maysir test --- yields a coherent verdict: premium-only
DEX perpetuals used on productive underlying assets with disciplined risk
management do not violate the three canonical Islamic prohibitions. The
\textit{riba} is absent by design. The \textit{gharar} is substantially reduced
by on-chain transparency. The \textit{maysir} determination is use-contingent,
as it is for all derivatives.

The practical implication is significant. 1.8 billion Muslims face structural
barriers to participation in the dominant leveraged instrument of modern crypto
markets. This paper demonstrates that those barriers are not inherent to the
instrument class --- they are a consequence of specific funding formula design
choices. Where the interest parameter is set to zero, as in dYdX~v4 and D8X,
the structural basis for prohibition is removed. This is a replicable design
criterion, not a platform-specific accident: any perpetual futures exchange that
adopts a premium-only funding formula satisfies the same structural condition.

The reform implication is therefore broader than a verdict on existing platforms.
It is a specification: perpetual futures can be structured for Shariah compliance
by design, and the empirical evidence presented here demonstrates that such
structures function at institutional scale across a full market cycle. We call on
Islamic scholars, Shariah boards, and standard-setting bodies --- particularly
AAOIFI and OIC Fiqh Academy --- to engage with the technical specifics of DEX
funding formulas and to develop formal guidance that enables institutional
participants to evaluate perpetual futures instruments on a principled,
formula-specific basis. The analytical framework is established; the empirical
validation is documented; the formal \textit{fatwa} must follow.

\textit{And Allah knows best.}

%% ============================================================
\sloppy
\begin{thebibliography}{99}

\bibitem{Ackerer2024}
Ackerer, D., Hugonnier, J., \& Jermann, U. (2024).
\textit{Perpetual Futures Pricing}.
NBER Working Paper 32936.
\url{https://www.nber.org/papers/w32936}

\bibitem{Alias2015}
Alias, M. (2015).
Gharar in forward and futures contracts.
\textit{ResearchGate}.
\url{https://www.researchgate.net/publication/273904700}

\bibitem{Alias2021}
Alias, M. (2021).
Derivative instruments: Criteria for speculation in Islamic finance.
\textit{Academy of Accounting and Financial Studies Journal}, 25(S3).
\url{https://www.abacademies.org/articles/derivative-instruments-criteria-for-speculation-in-islamic-finance-11080.html}

\bibitem{Alias2023}
Alias, M. (2023).
Derivative instruments: Criteria for speculation in Islamic finance (updated).
\textit{ResearchGate}.
\url{https://www.researchgate.net/publication/369559261}

\bibitem{AliasOthman2024}
Alias, M., \& Othman, A. (2024).
A review on \textit{gharar} dimension in modern Islamic finance transactions.
\textit{Journal of Islamic Accounting and Business Research}, 16(5), 976--994.
\url{https://www.emerald.com/insight/content/doi/10.1108/JIABR-02-2023-0060}

\bibitem{Amin2024}
Amin, M. (2024).
Are perpetual futures contracts Shariah compliant?
\url{https://www.mohammedamin.com/Islamic_finance/Perpetual-futures.html}

\bibitem{Ayub2007}
Ayub, M. (2007).
\textit{Understanding Islamic Finance}.
John Wiley \& Sons.

\bibitem{Binance2019}
Binance. (2019).
Introduction to Binance Futures funding rates.
\textit{Binance Support}.
\url{https://www.binance.com/en/support/faq/detail/360033525031}

\bibitem{AAOIFI2015}
AAOIFI. (2015).
\textit{Shari'ah Standards for Islamic Financial Institutions}.
Accounting and Auditing Organisation for Islamic Financial Institutions.
\url{https://www.iefpedia.com/english/wp-content/uploads/2017/12/Shariaa-Standards-ENG.pdf}

\bibitem{Chainalysis2025}
Chainalysis. (2025).
\textit{An Introduction to Perpetual Futures}.
\url{https://www.chainalysis.com/blog/perpetual-futures-intro-2025/}

\bibitem{CoinGecko2025}
CoinGecko. (2025).
\textit{State of Crypto Perpetuals Market 2025}.
\url{https://assets.coingecko.com/reports/2025/CoinGecko-State-of-Crypto-Perpetuals-Market.pdf}

\bibitem{D8X2024}
D8X Exchange. (2024).
The silent difference between DeFi perpetual futures contracts.
\textit{Medium}.
\url{https://medium.com/@d8x.exchange/the-silent-difference-between-defi-perpetual-futures-contracts-283e47a84b94}

\bibitem{DSNMUI2021}
DSN-MUI. (2021).
Fatwa No.\ 144/DSN-MUI/XI/2021 on crypto assets.
\textit{Indonesian Ulema Council}.

\bibitem{dYdX2024}
dYdX. (2024).
Perpetual funding rate.
\textit{dYdX Help Center}.
\url{https://help.dydx.exchange/en/articles/4797443-perpetual-funding-rate}

\bibitem{ElGamal2001}
El-Gamal, M. A. (2001).
An economic explication of the prohibition of \textit{gharar}.
\textit{Islamic Economic Studies}, 8(2), 29--58.

\bibitem{ElGamal2006}
El-Gamal, M. A. (2006).
\textit{Islamic Finance: Law, Economics, and Practice}.
Cambridge University Press.

\bibitem{Elhessi2018}
Elhessi, A. (2018).
\textit{Gharar} in futures contracts from an Islamic perspective.
\textit{Risk Governance and Control: Financial Markets and Institutions}.
\url{https://virtusinterpress.org/Gharar-in-futures-contracts-from-Islamic-perspective-A-case-study-of-USA-Malaysia.html}

\bibitem{Hayes2016}
Hayes, A. (2016).
Announcing the launch of the perpetual XBTUSD leveraged swap.
\textit{BitMEX Blog}.
\url{https://www.bitmex.com/blog/announcing-the-launch-of-the-perpetual-xbtusd-leveraged-swap}

\bibitem{Hyperliquid2025}
Hyperliquid. (2025).
Funding rate mechanics.
\textit{Hyperliquid Documentation}.
\url{https://hyperliquid.gitbook.io/hyperliquid-docs/trading/funding}

\bibitem{Indonesian2024}
Zulkarnain, R., \& Syibly, M. R. (2024).
Analysis of futures trading on the Binance platform from the perspective of Fiqh Muamalah.
\textit{Journal of Islamic Economics Lariba}.
\url{https://journal.uii.ac.id/JIELariba/article/download/33018/17999/132830}

\bibitem{IslamQA2023}
IslamQA. (2023).
Is it permissible to study futures contracts?
\url{https://islamqa.info/en/answers/241297}

\bibitem{Jobst2007}
Jobst, A. A. (2007).
Derivatives in Islamic finance.
\textit{SSRN Working Paper}.
\url{https://papers.ssrn.com/sol3/papers.cfm?abstract_id=1015615}

\bibitem{Kalimullina2020}
Kalimullina, M. (2020).
Islamic finance and food security: A review of the literature.
\textit{Heliyon}, 6(11).
\url{https://pmc.ncbi.nlm.nih.gov/articles/PMC7644896}

\bibitem{Kamali2000}
Kamali, M. H. (2000).
\textit{Islamic Commercial Law: An Analysis of Futures and Options}.
Islamic Texts Society.

\bibitem{Kasri2016}
Kasri, N. S. (2016).
\textit{Critical Analysis of Malaysian Securities Commission SAC Resolution on
Crude Palm Oil Futures}.
PhD Thesis, University of Gloucestershire.
\url{https://eprints.glos.ac.uk/3444/}

\bibitem{Malkawi2014}
Malkawi, B. H. (2014).
Financial derivatives between Western legal tradition and Islamic finance.
\textit{Journal of Banking Regulation}, 15(1).

\bibitem{Musaffa2024}
Musaffa Academy. (2024).
Things about perpetual futures you need to know.
\url{https://academy.musaffa.com/things-about-perpetual-futures-you-need-to-know/}

\bibitem{Obaidullah2005}
Obaidullah, M. (2005).
\textit{Islamic Financial Services}.
King Abdulaziz University.
\url{https://www.kau.edu.sa/Files/320/Researches/50893_21007.pdf}

\bibitem{PracticalIF2023}
Practical Islamic Finance. (2023).
Perpetual futures: Halal or haram?
\url{https://blog.practicalislamicfinance.com/perpetual-futures-halal-or-haram/}

\bibitem{ResearchGate2026}
Unpublished/Preprint. (2026).
Islamic sale and purchase principles in decentralised exchanges.
\textit{ResearchGate Preprint}.
\url{https://www.researchgate.net/publication/399063294}

\bibitem{Salamon2015}
Salamon, H. (2015).
Speculation: The Islamic perspective; A study on al-\textit{maisir}.
\textit{Mediterranean Journal of Social Sciences}, 6(4).
\url{https://www.academia.edu/25656159/}

\bibitem{SeekersGuidance2025}
SeekersGuidance. (2025, October 2).
Are crypto perpetual contracts still haram if there's no interest, and the
asset price tracks?
Answered by Shaykh Muhammad Carr; approved by Shaykh Faraz Rabbani.
\url{https://seekersguidance.org/answers/transactions/are-crypto-perpetual-contracts-still-haram-if-theres-no-interest-and-the-asset-price-tracks}

\bibitem{Shiller1993}
Shiller, R. J. (1993).
\textit{Macro Markets: Creating Institutions for Managing Society's Largest
Economic Risks}.
Oxford University Press.

\bibitem{SoleJobst2012}
Sol\'{e}, J., \& Jobst, A. A. (2012).
Operative principles of Islamic derivatives: Towards a coherent theory.
\textit{IMF Working Paper} WP/12/63.
\url{https://www.elibrary.imf.org/view/journals/001/2012/063/article-A001-en.xml}

\bibitem{Tashfeen2022}
Tashfeen, M. (2022).
Financial derivatives hedging: A review of Islamic finance at crossroads.
\textit{Global Economics Review}, 7(2).
\url{https://www.gerjournal.com/article/financial-derivatives-hedging-a-review-of-islamic-finance-at-cross-roads}

\bibitem{Uddin2015}
Uddin, M. A. (2015).
Principles of Islamic finance: Prohibition of \textit{riba}, \textit{gharar}
and \textit{maysir}.
\textit{MPRA Paper} No.~67711.
\url{https://mpra.ub.uni-muenchen.de/67711/1/MPRA_paper_67711.pdf}

\bibitem{Uluyol2024}
Uluyol, B. (2024).
Financial derivative instruments and their applications in Islamic banking
and finance.
\textit{Borsa Istanbul Review}, 24(2).
\url{https://www.sciencedirect.com/science/article/pii/S2214845024000334}

\bibitem{Ghazali2025}
Ghazali, N.~M., et al.\ (2025).
Exploring financial derivatives through Shariah principles.
\textit{Al-Qanatir: International Journal of Islamic Studies}.

\bibitem{Alsmadi2025}
Alsmadi, A.~A., et al.\ (2025).
Beyond compliance: Islamic Fintech and CSR.
\textit{Future Business Journal}, 11(1).

\bibitem{Farhat2025}
Farhat, A. (2025).
The impact of Shariah compliance on firm's performance and risk.
\textit{Review of Financial Economics}.

\bibitem{Modjo2025}
Modjo, M.~I., et al.\ (2025).
Weathering the storm: Shariah compliance, digital innovation, and stock
performance.
\textit{Journal of Islamic Monetary Economics and Finance}, 11(1).

\bibitem{SP2025}
S\&P Global. (2025).
\textit{Islamic Finance 2025--2026: Resilient Growth Amid Upcoming Headwinds}.
S\&P Global Ratings.
\url{https://www.spglobal.com/ratings/en/research/articles/250417-islamic-finance-2025-2026}

\end{thebibliography}

\end{document}
